\section{\bi{A different global minimizer at a single width and floor}{單一寬度與下限處的另一種全域極小者}}
\begin{thm}\label{th:point}
\bi{At $N=6$, $s=22$, $h=1/1000$, every minimizing gap word is a cyclic shift of $(h,h,R-h,1/2-R,1/2-R,R-h)$. A representative cell and strict enclosures are}{當 $N=6$、$s=22$、$h=1/1000$，每個極小間隙字均為 $(h,h,R-h,1/2-R,1/2-R,R-h)$ 的循環移位。一組代表與嚴格包圍為}
\begin{align}
 X_*&=(-h,0,h,R,\tfrac12,1-R)\pmod 1,\label{eq:point}\\
 0.249999032167246504146498709394&<R<0.249999032167246504146498709395,\notag\\
 0.66503584451339086827047430&<M_6(1/1000;22)<0.66503584451339086827047431 .\notag
\end{align}
\bi{There is one geometry and two species-preserving orbits, distinguished by the net charge of the triple, even if reflection is allowed.}{幾何形狀唯一；保留電荷種類時，由緊三簇的淨電荷區分為兩個軌道，即使容許反射亦然。}
\end{thm}
\begin{proof}
\bi{\emph{Comparison with the packed face.} Set $G=G_{22}$. Moving the third coordinate of the collision trial in \cite[Eq.~(17)]{infimum} from $125/1000$ to $126/1000$ gives the admissible cell $X=(0,125,126,250,500,750)/1000$. The signed weight of a pair of packed triples with directed center separation $d$ is $2A_n^2(1-\cos(2\pi nd))$. Contracting each triple to its net charge sets $A_n=1$ and gives auxiliary energy $E_0(d)=1+[G(d)-G(0)]/3$. The first nine terms of $E(Y)-E_0$ are nonnegative since $|A_n|\le1$. From $q<3/4$, $c<1$, $A_n^2-1\le8$, the loss in the tail is bounded by $(32/3)(3/4)^{100}/[1-(3/4)^{21}]$. Write $r=\operatorname{dist}(d,\mathbb Z)$. For $r\le1/4$ the nearest image alone exceeds $2e^{-11/2}$; otherwise convexity of $e^{-22u^2}$ on $[1/4,3/4]$ gives the same lower bound from the nearest two. Also $G(0)\le1+2e^{-22}/(1-e^{-66})$. Rational evaluation gives $E(Y)>.669391180769582$ for every $Y$ on the face and $E(X)<.667106810407357$. Thus $X$ beats the whole face by more than $.002284370362225$. The admissible $Z=(0,1,2,251,501,751)/1000$ already has $E(X)-E(Z)>.00207096588405$.}{\emph{能量上界與整族比較。} 記 $G=G_{22}$。將文獻 \cite[式 (17)]{infimum} 碰撞試探的第三個坐標由 $125/1000$ 移至 $126/1000$，得可行構型 $X=(0,125,126,250,500,750)/1000$。兩緊三簇的中點有向間距為 $d$ 時，帶號模態權重為 $2A_n^2(1-\cos(2\pi nd))$。將各簇收縮至淨電荷相當於置 $A_n=1$，輔助能量為 $E_0(d)=1+[G(d)-G(0)]/3$。$E(Y)-E_0$ 的前九項因 $|A_n|\le1$ 皆非負。由 $q<3/4$、$c<1$、$A_n^2-1\le8$，其餘尾項的損失不逾 $(32/3)(3/4)^{100}/[1-(3/4)^{21}]$。設 $r=\operatorname{dist}(d,\mathbb Z)$。若 $r\le1/4$，單一最近像已超過 $2e^{-11/2}$；否則由 $e^{-22u^2}$ 在 $[1/4,3/4]$ 的凸性，最近二像給出同一下界。又 $G(0)\le1+2e^{-22}/(1-e^{-66})$。有理估計得整面 $E(Y)>.669391180769582$、$E(X)<.667106810407357$。故 $X$ 比整個面的任一構型低逾 $.002284370362225$。可行點 $Z=(0,1,2,251,501,751)/1000$ 已有 $E(X)-E(Z)>.00207096588405$。}

\bi{\emph{Global localization.} Cutting at all slack gaps gives one of the twelve cyclic block types in Table~\ref{tab:point}. These exhaust the nonempty subsets of six cuts modulo rotation and reversal, since $6h<1$. For two disjoint contacting pairs, remove them and let $r=1/2-\delta$ be the short distance of the residual opposite pair. The fixed-gauge auxiliary energy is $E_0(r)$. With $K=44$, $\|G''\|_\infty=-G''(0)<K$ by positive Fourier coefficients and the positive noncentral image contributions to $G''(0)$. The two pair self terms and their mixed double difference add at most $4Kh^2$ to the ordered score. On either residual arc of length $u$, the interaction of a removed pair with the residual pair is a first difference of $G(v)-G(u-v)$; its derivative is $G'(v)+G'(u-v)$. The kernel is decreasing up to $1/2$. At most one argument exceeds $1/2$, by at most $\delta$; $G'(1/2)=0$ bounds this derivative by $K\delta$. Adding two pairs in both directions costs at most $4Kh\delta$. The certified signs $G''<0$ on $[0,.1]$, $G'<0$ on $[.1,.25]$, $G''>1$ on $[.25,.5]$ establish monotonicity. They and $G(.25)-G(.5)>1/8$ give $G(r)-G(.5)\ge\delta^2/2$ by two integrations or monotonicity. Completing the square gives}{\emph{全域定位。} 在所有鬆弛間隙切開圓周，得到表~\ref{tab:point} 的十二種循環塊型；因 $6h<1$，它們窮盡六個切口的非空子集在旋轉及反序下的代表。若有兩對不共享粒子的接觸，移去它們後，餘下的異號粒子短弧距離設為 $r=1/2-\delta$，其輔助能量為 $E_0(r)$。取 $K=44$，由 Fourier 係數為正及 $G''(0)$ 非中央像的正貢獻知 $\|G''\|_\infty=-G''(0)<K$。兩對的自身項及其二重差分對有向得分共增加至多 $4Kh^2$。在長 $u$ 的餘下弧上，被移去一對與兩餘點的交互項是 $G(v)-G(u-v)$ 的一階差分，其導數為 $G'(v)+G'(u-v)$。核至 $1/2$ 遞減；兩自變量至多一個超過 $1/2$，且超出至多 $\delta$，由 $G'(1/2)=0$ 得導數至多 $K\delta$。兩對及雙向計數合計至多 $4Kh\delta$。經認證的 $G''<0$ 於 $[0,.1]$、$G'<0$ 於 $[.1,.25]$、$G''>1$ 於 $[.25,.5]$ 給出單調性。以兩次積分或單調性配合 $G(.25)-G(.5)>1/8$ 得 $G(r)-G(.5)\ge\delta^2/2$。完成平方即得}
\[
 E\ge E_0(1/2)+\delta^2/6-2Kh\delta/3-2Kh^2/3
 \ge E_0(1/2)-(2K^2+2K)h^2/3>.66807 .
\]
\bi{If all six gaps are slack, a single-site variation of the ordered score requires $\sigma_i V''(x_i)\le G''(0)$, $V(u)=\sum_j\sigma_j G(u-x_j)$. Truncate $V''$ at degree two to a trigonometric polynomial $T$. Its tail divided by $2c(2\pi)^2$ is at most $54(.64)^9/[1-(4/3)^2(.64)^7]<.63+4(.63)^4$, while $-G''(0)/[2c(2\pi)^2]\ge q+4q^4$, $.63<q<.64$. Thus $\sigma_i T(x_i)<0$ for all six sites. Alternation forces six distinct zeros; multiplication by $e^{4\pi i u}$ gives a nonzero polynomial of degree at most four in $e^{2\pi i u}$. This is impossible.}{若六隙全鬆弛，由有向得分的單粒子變分有 $\sigma_i V''(x_i)\le G''(0)$，其中 $V(u)=\sum_j\sigma_j G(u-x_j)$。截取 $V''$ 至二階三角多項式 $T$。尾項除以 $2c(2\pi)^2$ 後至多為 $54(.64)^9/[1-(4/3)^2(.64)^7]<.63+4(.63)^4$；而 $-G''(0)/[2c(2\pi)^2]\ge q+4q^4$、$.63<q<.64$。所以全部六點皆有 $\sigma_iT(x_i)<0$。符號交替迫使六個相異零點；但乘 $e^{4\pi iu}$ 後得到關於 $e^{2\pi iu}$ 的非零至多四次多項式，矛盾。}
\par
\begin{table}[tbp]\centering
\caption{\bi{Exhaustion at $(22,1/1000)$. Digits count sites in maximal packed blocks.}{參數點 $(22,1/1000)$ 的窮盡分類。每個數字是一個極大緊塊的粒子數。}}\label{tab:point}
\begin{tabular}{lp{.52\linewidth}}\toprule
\bi{Cyclic words}{循環塊長字} & \bi{Certificate}{排除或定位證書}\\\midrule
$6,33,42,51,222,321,411,2121,2211$ & $E>.66807>E(X)$\\
$111111$ & \bi{six zeros of a degree-two trigonometric polynomial}{二次三角多項式須有六個零點}\\
$21111$ & \bi{1,269 boxes, no survivors}{1,269 盒，全部排除}\\
$3111$ & \bi{541 boxes, 10 leaves in the uniqueness region}{541 盒，10 個葉盒進入唯一性區}\\\bottomrule
\end{tabular}\end{table}
\par
\bi{For the last two faces, the block-kernel inequalities of Appendix~\ref{sec:cert} exclude a box if $W<W(X)$ or a two-sided translation violates necessary gradient or curvature conditions. In the single-contact tree the three exclusion counts are $36,283,316$; in the adjacent-contact tree they are $24,158,79$, and the remaining ten closed boxes have all four slack gaps in $[.21,.29]$. No boundary is lost because any newly contacting gap has its own actual block word. A second full cover in block-start coordinates visits 5,638 boxes with $344,32,934,1509$ leaves discarded by order, score, gradient and curvature, and six localized leaves, all of type $3111$. Both covers use $X$ as the feasible bound.}{最後兩面採附錄~\ref{sec:cert} 的塊核估計，若 $W<W(X)$ 或雙向平移違反梯度或曲率必要條件即排除整盒。單接觸樹的三類排除葉數為 $36,283,316$；相鄰雙接觸樹為 $24,158,79$，另十個閉盒內四個鬆弛間隙皆屬 $[.21,.29]$。任何新接觸的邊界點另歸其實際塊長字。另以塊首坐標全域覆蓋 5,638 盒，次序、得分、梯度、曲率分別排除 $344,32,934,1509$ 個葉盒，六個保留葉盒均為 $3111$。兩種覆蓋都以 $X$ 作可行界。}

\bi{\emph{Local shape and equality.} With a $+,-,+$ triple write its coordinates $(-v,u-v,u)$, fixing its signed first moment. The other sites $r,z,t$ have signs $\tau=(-,+,-)$. For short gaps $h\le u,v\le.02$ and the four long gaps in $[.21,.29]$, differentiation of the score gives}{\emph{局部形狀與等號。} 將 $+,-,+$ 三點寫為 $(-v,u-v,u)$，固定其帶號一階矩。其餘點 $r,z,t$ 帶號 $\tau=(-,+,-)$。當短隙 $h\le u,v\le.02$ 且四個長隙在 $[.21,.29]$，對得分微分有}
\[
 W_u/v=\int_0^1\left[G''(u+\lambda v)+
             \sum_{y=r,z,t}\tau_yG''(y-u+\lambda v)\right]\dd\lambda .
\]
\bi{For $W_v/u$ replace the arguments by $v+\lambda u$ and $y+v-\lambda u$. Verified bounds on the three argument ranges are $G''<-38$ on $[0,.04]$, $G''>14$ on $[.21,.33]$, $G''<7$ on $[.42,.58]$. Both derivatives divided by the other gap are $<-59$. Contraction to $u=v=h$ shrinks the external argument intervals and is feasible. Set $C(d)=G(d-h)-G(d)+G(d+h)$. The variable score becomes}{對 $W_v/u$ 將兩個自變量換成 $v+\lambda u$ 與 $y+v-\lambda u$。三段核界為 $G''<-38$ 於 $[0,.04]$、$G''>14$ 於 $[.21,.33]$、$G''<7$ 於 $[.42,.58]$。兩偏導數各除以另一短隙均 $<-59$。收縮至 $u=v=h$ 會縮小外點自變量區間，並保持可行。令 $C(d)=G(d-h)-G(d)+G(d+h)$，可變得分為}
\[
 -C(r)+C(z)-C(t)-G(z-r)+G(t-r)-G(t-z).
\]
\bi{In the convex set where the four consecutive distances of $0,r,z,t$ are in $[.21,.31]$, $C''>14$ on the side distances and $C''<7$ on the diagonals, with the same bounds for $G''$. In the Hessian the four side coefficients are strictly below $-14$, the two diagonal coefficients below 7. Each squared diagonal velocity difference is bounded by the sum of all four squared side differences, averaging the two two-edge paths and applying $(a+b)^2\le2(a^2+b^2)$. Any nonzero velocity after fixing translation has a nonzero side difference. Strict concavity follows. Reflection sends $(r,z,t)$ to $(1-t,1-z,1-r)$, so at a point on its fixed line the differential equals the differential on the averaged, tangent direction. On that line stationarity reads}{在 $0,r,z,t$ 的四個相鄰距離屬於 $[.21,.31]$ 的凸域上，邊距有 $C''>14$，對角距有 $C''<7$，$G''$ 滿足相同界。Hessian 中四邊係數嚴格小於 $-14$，兩對角係數小於 7。對於一條對角，沿兩條二邊路徑分別用 $(a+b)^2\le2(a^2+b^2)$ 再取平均，其速度差平方至多為四邊速度差平方之和。固定平移後任一非零速度至少使一邊有非零速度差，遂得嚴格凹性。反射將 $(r,z,t)$ 送至 $(1-t,1-z,1-r)$，因此在不動直線上一階微分等於其對平均後切向的微分。直線上的駐點條件為}
\begin{equation}\label{eq:force}
 C'(R)-G'(1/2-R)-G'(2R)=0 .
\end{equation}
\bi{The image-sum enclosures give strictly negative and positive signs at the left and right endpoints in \eqref{eq:point}. The root in the convex domain is a full stationary point and strict concavity makes it unique. Compactness supplies a global maximizer and the exhaustive localization puts it in this domain. Substitution of the entire root interval into the fifteen pair terms proves the energy enclosure. In particular $E(X)-M_6>.00207096589396576$. Reflection preserves the charge of the uniquely identifiable triple, whereas global reversal interchanges its two assignments.}{像和包圍在式 \eqref{eq:point} 的左右端點分別給出嚴格負、正號。凸域內的根是完整駐點，且由嚴格凹性唯一。緊性保證全域極大者存在，而窮盡定位迫使它落入此域。將根的整個包圍區間代入十五項配對能量得所述嚴格能量界，特別有 $E(X)-M_6>.00207096589396576$。反射保留唯一可辨緊三簇的淨電荷，整體換號則互換兩種指派。}
\end{proof}

